\section{forme differenziali lineari}

Se $f\colon A\subset \RR^n \to \RR$ è una funzione differenziabile il suo 
differenziale $df \colon A \to \mathcal L (\RR^n, \RR)$ è una forma differenziale lineare
($1$-forma).
Si tratta di una funzione che ad ogni punto $x\in A$ associa una applicazione lineare 
$df(x)\colon \RR^n \to \RR$ e dunque ad ogni $v\in \RR^n$ associa un numero reale 
$df(x)[v]$. 
Se sfruttiamo il prodotto scalare standard su $\RR^n$ 
possiamo identificare l'applicazione lineare $df(x)$ con il vettore $\nabla f(x)$:
\[
    df(x)[v] = \frac{\partial f}{\partial v}(x) = \nabla f(x) \cdot v.
\]
Infatti il prodotto scalare su uno spazio vettoriale $V$ 
induce un isomorfismo $V \to V^*$, $v \mapsto (w \mapsto v \cdot w)$.

\begin{definition}[forme differenziali lineari]
    Sia $A\subset \RR^n$ un aperto.
    Denotiamo con $\Omega^1(A)$ l'insieme di tutte le funzioni $\omega\colon A \to (\RR^n)^*$
    che ad ogni punto $x\in A$ associano una applicazione lineare $\omega(x)\colon \RR^n \to \RR$,
    $v\mapsto \omega(x)[v]$.
\end{definition}

\begin{remark}
Le forme differenziali sono funzioni di due variabili: $(x,v)\mapsto \omega(x)[v]$.
La dipendenza da $x$ non è lineare (e in generale lo spazio ambiente potrebbe non essere uno 
spazio vettoriale, dunque non avrebbe senso richiedere la linearità in $x$)
mentre la dipendenza da $v$ è lineare. 
A volte si fissa una delle due variabili ma non l'altra.
Per questo utilizziamo le parentesi quadre per distinguere la seconda variabile dalla prima.
In generale useremo le parentesi quadre per ricordare che la dipendenza da quella variabile è lineare.

Le forme differenziali si definiscono più in generale sulle \emph{varietà differenziabili}, che sono una 
generalizzazione a più dimensioni delle superfici regolari.
La variabile $x$ rappresenta un punto sulla superficie mentre $v$ è un vettore tangente 
a quel punto. Lo spazio prodotto dove vive la coppia $(x,v)$ si chiama \emph{fibrato tangente}.

Purtroppo sullo spazio piatto $\RR^n$ si può facilmente fare confusione tra le due varibili 
perché lo spazio tangente ad $\RR^n$ in ogni punto è $\RR^n$ stesso.
Bisogna pensare ad $x\in \RR^n$ come un punto, 
mentre $v\in \RR^n$ è un vettore applicato in quel punto.
\end{remark}

Sappiamo, dall'algebra lineare, che il duale di uno spazio vettoriale 
di dimensione finita ha la stessa dimensione dello spazio stesso. 
Inoltre se $e_1,\dots, e_n$ è la base canonica di $\RR^n$,
la corrispondente base dello spazio duale $(\RR^n)^*$ è data 
dalle applicazioni lineari $\lambda_1,\dots, \lambda_n$ definite 
da 
\[
  \lambda_k(e_j)=\delta_{kj} = \begin{cases}
  1 & \text{se } k=j,\\
  0 & \text{altrimenti.}
  \end{cases}
\]
Se denotiamo con $x_k\colon \RR^n \to \RR$ la funzione $\vec x \mapsto x_k$ 
che restituisce la $k$-esima coordinata di un punto $\vec x = (x_1,\dots, x_n)$, 
possiamo osservare che $dx_k[e_j] = \frac{\partial x_k}{\partial x_j} = \delta_{kj}$ 
e dunque $dx_k = \lambda_k$ è in effetti la base canonica di $(\RR^n)^*$. 
Significa che ogni $\omega \in \Omega^1(A)$ può essere scritto 
nella forma 
\[
   \omega(p) = \xi_1 dx_1 + \xi_2 dx_2 + \dots + \xi_n dx_n.
\]
Dunque c'è un isomorfismo $\omega \mapsto \xi$ che mette in 
corrispondenza le forme differenziali lineari $\omega \in \Omega^1(A)$ 
con i campi vettoriali $\xi \colon A\to \RR^n$.
In questo isomorfismo:
\[
  \omega(p)[v] = \xi(p) \cdot v
\]
(abbiamo qui sfruttato il fatto che sullo spazio $\RR^n$ è fissato un prodotto 
scalare canonico che induce un isomorfismo canonico tra $\RR^n$ 
e il suo duale).

\section{integrale di una forma differenziale}

Se $u\colon [a,b]\to A\subset \RR^n$ è una curva regolare e $\omega\in \Omega^1(A)$ 
è una $1$-forma differenziale, possiamo definire l'integrale della forma differenziale lungo la curva:
\[
 \int_u \omega \defeq \int_a^b \omega(u(t))[u'(t)]\, dt.
\]
Osserviamo che se $\omega$ si rappresenta con il campo vettoriale $\xi$, si ha $\omega(x)[v] = \xi(x)\cdot v$ 
e dunque 
\[
 \int_u \omega = \int_a^b \xi(u(t))\cdot u'(t)\, dt.
\]
Questo integrale si chiama \emph{circuitazione} di $\xi$ lungo $u$.
Osserviamo che questo integrale è indipendente dalla parametrizzazione della curva $u$
(se si mantiene l'orientazione).
Infatti se $\phi\colon [A,B]\to [a,b]$ con $\phi(A)=a$ e $\phi(B)=b$ 
posto $t=\phi(s)$ si ha $(u\circ \phi)'(s)\, ds = u'(\phi(s)) \phi'(s)\, ds = u'(t)\, dt$.
Dunque la formula del cambio di variabile garantisce l'uguaglianza
\begin{align*}
\int_{u\circ \phi} \omega 
&= \int_A^B \omega(u(\phi(s)))[(u\circ \phi)'(s)]\, ds
= \int_A^B \omega(u(\phi(s)))[u'(\phi(s))] \phi'(s)\, ds\\
&= \int_a^b \omega(t)[u'(t)]\, dt = \int_u \omega.
\end{align*}

In generale se $\phi \colon A\subset \RR^m \to \RR^n$ 
il differenziale $d\phi(x)$ calcolato in un punto $x\in A$
manda i vettori di $\RR^m$ nei vettori di $\RR^n$ 
e permette quindi di definire una forma differenziale su $A$
se abbiamo una forma differenziale su $\phi(A)$
(la forma differenziale viene \emph{tirata indietro} dal codominio al dominio di $\phi$).
Formalmente si definisce il \emph{pull-back}
$\phi^*\colon \Omega^1(\phi(A)) \to \Omega^1(A)$ 
tramite la formula
\[
   (\phi^* \omega)(x)[v] = \omega(\phi(x))[d\phi(x)[v]].
\]

Osserviamo che se $u\colon I\to A\subset \RR^n$ 
con $I$ un intervallo di $\RR$ 
allora $\phi^* \omega \in \Omega^1(I)$. 
Ma questo spazio ha dimensione $1$, ed è generato dal differenziale 
$dt$ dell'identità $t\mapsto t$. 
In effetti se pensiamo all'integrale 
$\int_a^b f(t)\, dt$
come all'integrazione della forma differenziale $f(t)\, dt$ 
sull'intervallo $[a,b]$, possiamo notare che si ha
\[
 \int_u \omega = \int_I u^* \omega
\]
in quanto (qui $e_1=1$ è l'unico vettore della base canonica di $\RR$)
\[
  u^* \omega(t)[e_1] = \omega(u(t))[u'(t)]\, dt.
\]

\section{forme chiuse, forme esatte}

\begin{definition}[forma esatta]
Una forma differenziale $\omega\in \Omega^1(A)$ si dice essere \emph{esatta}
se esiste $f\colon A\to \RR$ tale che $\omega = df$.
\end{definition}

L'integrale di un differenziale esatto risulta banale. Se $\omega = df$
\begin{align*}    
    \int_u \omega 
    &= \int_u df 
    = \int_a^b df(u(t))[u'(t)]\, dt 
    = \int_a^b \frac{\partial f(u(t))}{u'(t)}\, u'(t) \, dt\\
    &= \int_a^b (f\circ u)'(t) \, dt 
    = f(u(b)) - f(u(a)).
\end{align*}
In questo caso se due curve hanno gli stessi estremi l'integrale è lo stesso,
e se la curva è chiusa, cioè se $u(b) = u(a)$,
l'integrale è nullo. 
Se dobbiamo calcolare l'integrale di una forma differenziale $\omega$ sarà dunque 
interessante capire se tale forma è esatta.

Osserviamo innanzitutto che non tutte le forme differenziali lineari 
sono il differenziale di una funzione $f\colon A\to \RR$.
Infatti se $f\in C^2(A)$ e $\omega = df$
allora 
\[
    \omega = \xi_1 dx_1 + \dots + \xi_n dx_n
\]
e, per il teorema di Schwartz si ha 
\begin{equation}\label{eq:forma-chiusa}
  \frac{\partial \xi_k}{\partial x_j} = \frac{\partial \xi_j}{\partial x_k}.
\end{equation}

\begin{definition}[forma chiusa]
Sia $\omega\in \Omega^1(A)$ una $1$-forma di classe $C^1$
ovvero tale che esiste $\xi\in C^1(A,\RR^n)$ tale che 
$\omega(x)[v] = \xi(x)\cdot v$. 
Diremo che $\omega$ è chiusa se vale~\eqref{eq:forma-chiusa}
\end{definition}

Se $df = \omega$ diremo che $f$ è una primitiva di $\omega$ 
e diremo che $\omega$ è integrabile. 
Abbiamo quindi dimostrato il seguente teorema.

\begin{theorem}[condizione necessaria per l'integrabilità]
\label{th:forma-chiusa}
Se $\omega$ è una $1$-forma di classe $C^1$ ed esiste 
$f\in C^1(A)$ tale che $\omega = df$ allora $\omega$ è chiusa. 
\end{theorem}

Per le funzioni di una variabile la codizione~\eqref{eq:forma-chiusa} è vuota, 
perché necessariamente $k=j=1$. 
E, in effetti, tutte le funzioni continue, in una variabile, sono integrabili 
per il teorema fondamentale del calcolo. 
In più variabili la condizione non è vuota. 

\begin{example}[forma differenziale non chiusa]
Sia $\omega(x,y) = y dx$, $\omega \in \Omega^1(\RR^2)$ è una forma di classe $C^\infty$ 
associata al campo vettoriale $\xi(x,y) = (y,0)$.
L'equazione~\eqref{eq:forma-chiusa} diventa $1=0$ e dunque questa forma non è chiusa e non 
può essere il differenziale di una funzione.
\end{example}

\begin{theorem}[Lemma di Poincaré]
    \label{th:lemma-poincare}
Sia $A$ un aperto \emph{stellato} di $\RR^n$
(cioè esiste $x\in A$ tale che per ogni $y\in A$ si ha $[x,y]\subset A$). 
Sia $\omega\in \Omega^1(A)$ una forma differenziale chiusa.
Allora $\omega$ è esatta 
(cioè esiste $f\colon A\to \RR$ tale che $df = \omega$).
\end{theorem}
%
\begin{proof}
Senza perdita di generalità possiamo supporre $x=0 \in A$.
Sia $\xi$ il campo vettoriale che rappresenta $\omega$.
Definiamo $f\colon A\to\RR$ come
\[
   f(x) = \int_{[0,x]} \omega = \int_0^1 \xi(tx)\cdot x\, dt 
   = \int_0^1 \sum_j x_j \xi_j(tx) \, dt.
\]
Ricordiamo che se $\omega$ è chiusa si ha $\frac{\partial \xi_j}{\partial x_k} = \frac{\xi_k}{\partial x_j}$.
Si ha 
\begin{align*}
 \frac{\partial}
 {\partial x_k}\sum_j x_j \xi_j(tx)
 &=  \sum_j \sum_i \delta_{ik} t x_j \frac{\partial \xi_j}{\partial x_i}(tx)
    + \sum_j \delta_{jk} \xi_j(tx)\\
 &= \sum_j t x_j \frac{\partial \xi_j}{\partial x_k}(tx) + \xi_k(tx)\\
 &= \sum_j t x_j \frac{\partial \xi_k}{\partial x_j}(tx) + \xi_k(tx)\\
 &= t \frac{d}{dt} \xi_k(tx) + \xi_k(tx) 
  = \frac{d}{dt} \enclose{t \xi_k(tx)}. 
\end{align*}
Dunque 
\[
 \frac{\partial f}{\partial x_k}(x) = \int_0^1 \frac{d}{dt} \enclose{t\xi_k(tx)}\, dt 
 = \Enclose{t\xi_k(tx)}_0^1 = \xi_k(x).
\]
Significa che $\omega = df$ come volevamo dimostrare.
\end{proof}


\begin{example}[forma chiusa ma non esatta]
Si consideri $A=\RR^2\setminus\ENCLOSE{(0,0)}$.
Sia $\omega\in \Omega^1(A)$ 
\[
\omega \defeq \frac{- y dx + x dy}{x^2+y^2}.
\]
Allora $\omega$ è chiusa ma non è esatta.
\end{example}
\begin{proof}
Poniamo
\[
a(x,y) = \frac{-y}{x^2+y^2},
\qquad 
b(x,y) = \frac{x}{x^2+y^2}
\]
cosicché $\omega = a(x,y)\,dx + b(x,y)\,dy$.
Allora da verifica diretta si osserva che 
\[
 \frac{\partial a}{\partial y} =
 \frac{\partial b}{\partial x} =
 \frac{y^2 - x^2}{(x^2+y^2)^2}
\]
e quindi effettivamente $\omega$ è chiusa.
Se la forma fosse esatta si avrebbe $\int_u \omega=0$ 
su ogni curva chiusa $u$.
Prendiamo invece $u\colon [0,2\pi]\to A$ la circonferenza
di raggio $\rho$:
$u(t) = \rho (\cos t,\sin t)$.
Cioè $x=\rho \cos t$, $y=\rho \sin t$, 
$dx = -\rho \sin t\,dt$, $dy=\rho \cos t\,dt$
da cui
\begin{align*}
\int_u \omega 
&= \int_0^{2\pi} \frac{-\rho \sin t (-\rho \sin t) + \rho \cos t \rho \cos t}{\rho^2\cos^2 t + \rho^2 \sin^2 t} \, dt\\
&= \int_0^{2\pi} 1\, dt = 2\pi. 
\end{align*}
Dunque $\omega$ non può essere esatta.
\end{proof}

L'esempio precedente non viola il teorema~\ref{th:lemma-poincare}
perché $A$ è un esempio di aperto non stellato.
Osserviamo che se da $A$ rimuoviamo l'asse delle $y$,
la funzione 
\[
  f(x,y) = \arctg \frac{y}{x}
\]
è ben definita e si ha
\[
df = \frac{1}{1+ \frac{y^2}{x^2}}\cdot 
    \enclose{\frac{-y}{x^2}\, dx + \frac{1}{x}\, dy}
    = \frac{-y\, dx + x\, dy}{x^2+y^2}.
\]
Si capisce dunque che l'esattezza di una forma chiusa dipende sostanzialmente 
dal dominio in cui la forma è definita.
